\section{Literature Review}

\textit{Classical exact methods.} Fractional LP bounds have a long history in discrete optimization \cite{ref1}. The Horowitz--Sahni meet-in-the-middle (MitM) algorithm splits the items into two halves and combines partial solutions, providing a classical $O(2^{n/2})$ arithmetic-operation baseline \cite{ref2}. Kellerer, Pferschy, and Pisinger consolidate dynamic programming, greedy bounds, and branch-and-bound (B\&B) \cite{ref3}. The \emph{core} concept was introduced by Balas and Zemel \cite{balas1980core}; Pisinger subsequently showed why core difficulty must be assessed across capacities and instance distributions \cite{ref4}. MINKNAP uses an expanding-core dynamic program \cite{pisinger1997minimal}, while COMBO combines dynamic programming with strong bounds \cite{martello1999combo}. These methods provide stronger classical context than brute-force baselines alone. Beier and V\"ocking's expected-polynomial knapsack analysis \cite{beier2004random} and Dey, Dubey, and Molinaro's random binary-IP B\&B analysis \cite{dey2021random} require their respective distributional assumptions; neither establishes polynomial search trees for the benchmark classes considered here.

\textit{Quantum search and its limits.} Grover's algorithm searches an unstructured space of size $N$ using $O(\sqrt{N})$ oracle queries \cite{ref8}. Amplitude amplification generalizes this principle \cite{ref11}, with optimization variants via pure adaptive search \cite{ref12} and D\"urr--H{\o}yer minimum finding \cite{ref13}. Searching all binary assignments therefore requires $O(2^{n/2})$ queries before accounting for oracle cost. Bennett et al.'s black-box lower bounds \cite{ref9} concern unstructured oracle access; they do not establish an unrestricted lower bound for structured knapsack algorithms or equate query complexity with runtime.

\textit{Subset-sum algorithms and memory assumptions.} Closely related subset-sum work includes the representation technique of Becker, Coron, and Joux \cite{ref5} and the quantum algorithms of Bernstein et al.\ \cite{ref16} and Helm--May \cite{ref17}. Their input models and, where used, heuristic assumptions must be distinguished from worst-case optimization with separate weights and values. Bonnetain et al.\ identify a problem with a walk-update heuristic and develop alternative algorithms and memory tradeoffs \cite{bonnetain2020improved}; this is not a rejection of every quantum subset-sum algorithm. Jaques and Rattew distinguish QRAM architectures and their resource assumptions \cite{ref18}. Such costs require explicit accounting, rather than a blanket assertion that quantum-addressable memory is impossible.

\textit{Gate-model knapsack and hybrid optimization.} Hou and Perkowski give gate-model bounded-knapsack circuits that test complete candidate assignments \cite{hou2020}. Hamitouche et al.\ combine Grover-based feasibility and threshold tests with a separately computed classical upper bound, for which a fractional-knapsack bound is one option \cite{hamitouche2025}. Wilkening et al.\ present a quantum tree generator with amplitude amplification; separately, their quantum B\&B benchmarks use fractional-greedy bounds \cite{wilkening2025}. Their released R1 software and experiment data provide a version-specific implementation reference \cite{wilkening2025code}. These works make quantum knapsack arithmetic and the use of fractional bounds prior art. Christiansen et al.\ combine feasible-solution state preparation with Grover-mixer QAOA \cite{ref23}, within the variational framework introduced by Farhi et al.\ \cite{ref25}. Christiansen et al.\ prove eventual optimal sampling at sufficiently large depth \cite{ref23}; this does not establish a practical finite-depth advantage. Grover adaptive search for constrained polynomial binary optimization \cite{ref26} is a distinct search approach, rather than a variational algorithm.

\textit{LP and Lagrangian bounds in quantum methods.} Sanavio et al.\ combine classical B\&B with quantum annealing for integer linear problems \cite{ref36}. Bo\.{z}ejko et al.\ use Lagrangian multiplier bounds in a D-Wave knapsack B\&B framework, including minimization over a multiplier set and integer rounding in pruning \cite{bozejko2026}. Their large-instance device experiments evaluate bounds at the root, rather than executing the full search tree on the QPU. Gabbassov et al.\ study Lagrangian duality for constrained quantum optimization and compare it with QUBO reformulations \cite{gabbassov2025}. Thus, neither Lagrangian relaxation nor taking an integer floor of a valid upper bound is itself a new contribution.

\textit{Query frameworks and reversible implementation.} Montanaro's quantum backtracking \cite{ref35} and quantum B\&B \cite{ref24} give query reductions under specified access to tree and bounding oracles. These frameworks do not make the cost of implementing an LP oracle free. Martiel and Remaud study concrete circuit implementation of quantum backtracking \cite{martiel2019}. Reversible arithmetic includes the networks of Vedral, Barenco, and Ekert \cite{vedral1996}, Draper's Fourier-basis addition \cite{draper2000}, and Cuccaro et al.'s ripple-carry addition \cite{cuccaro2004}; elementary controlled-gate decompositions are treated by Barenco et al.\ \cite{barenco1995}. Thapliyal et al.'s reversible integer divider \cite{thapliyal2016} is a useful arithmetic comparison, but the constant-coefficient oracle construction considered here does not execute that divider. Minimum interval stabbing is also an established combinatorial problem \cite{erickson2019algorithms}, so any contribution using it must concern the reduction and its endpoint conditions, rather than the greedy algorithm itself. Preskill's discussion of noise and depth constraints \cite{ref31} motivates separating logical circuit resources from device performance. Qiskit \cite{ref33} supplies software tools; its use alone is not evidence of execution on a quantum processor.

\textit{Closest prior work and our position.} The relevant distinction is between computing a fractional bound, testing a complete assignment, and reversibly deciding whether the LP relaxation of a partial prefix excludes an integer target. The construction under study targets the last task: for a fixed exact-efficiency prefix order, a classically certified bank of affine bounds is compiled into a reversible predicate that preserves the integer-threshold LP-pruning decisions. This scope requires positive integer weights, nonnegative integer profits, a coverage certificate, representable integer thresholds, a separate capacity check, and clean initial workspace. It does not imply reproduction of the numerical fractional LP value, support for arbitrary branching orders, or an implementation of the complete quantum walk. Circuit comparisons reconstructed from published fractional-greedy primitives \cite{wilkening2025,wilkening2025code} must therefore account for both predicate semantics and resource costs, and be identified as reconstructions rather than released LP-oracle circuits. Table~\ref{tab:related} summarizes the roles of the closest approaches without asserting historical priority or a hardware advantage.

\begin{table}[t]
\caption{Related methods and the distinctions relevant to an LP-pruning oracle. These descriptions identify scope, not a ranking of performance.}
\label{tab:related}
\begin{center}
\footnotesize
\setlength{\tabcolsep}{5pt}
\begin{tabular}{p{0.20\textwidth}p{0.30\textwidth}p{0.38\textwidth}}
\toprule
\textbf{Work} & \textbf{Computation or model} & \textbf{Relevant scope} \\
\midrule
Montanaro \cite{ref35,ref24} & Quantum backtracking and B\&B query frameworks & Assumes access to the required tree and bounding oracles. \\
Hou--Perkowski \cite{hou2020}; Hamitouche et al.\ \cite{hamitouche2025} & Gate-model candidate feasibility and objective tests & Complete-assignment tests; Hamitouche et al.'s optional fractional bound is classical. \\
Wilkening et al.\ \cite{wilkening2025,wilkening2025code} & QTG search; separate fractional-bound QBnB comparison & Published arithmetic primitives and released simulator/resource models; a matched LP-oracle circuit requires reconstruction. \\
Christiansen et al.\ \cite{ref23} & Feasible-solution preparation and QAOA & Variational search with depth-dependent behavior. \\
Sanavio et al.\ \cite{ref36}; Bo\.{z}ejko et al.\ \cite{bozejko2026} & Classical B\&B with quantum-annealing subproblems & A different computational model from a coherent reversible prefix predicate. \\
Martiel--Remaud \cite{martiel2019} & Circuit implementation of quantum backtracking & Implementation resources remain separate from abstract query counts. \\
\bottomrule
\end{tabular}
\end{center}
\end{table}

